\subsection{Finite Dimensional Inner Product Spaces}

\begin{definition}[Inner Product Space]
    Let $V$ be a vector space over $\mathbb{C}$ and $\langle-,-\rangle: V\times V\to\mathbb{C}$ be a function.
    We say the pair $(V,\langle-,-\rangle)$ is a complex inner product space if it satisfies:
    \begin{enumerate}
        \item (Linearity in 1st component) For any $u,v,w\in V$ and $\alpha,\beta\in\mathbb{C},$:
        $$\langle\alpha u+\beta v,w\rangle = \alpha\langle u, w\rangle + \beta\langle v,w\rangle.$$
        \item (Conjugate Symmetry) For any $u,v\in V$:
        $$\langle u,v\rangle = \overline{\langle v,u\rangle}.$$
        \item (Non-negativity) For any $u\in V$:
        $$\langle u,u\rangle\ge 0.$$
        \item (Non-degeneracy) For any $u\in V$:
        $$\norm{u}=\sqrt{\langle u,u\rangle}=0$$ 
        iff $u=0$.
    \end{enumerate}
    And $\langle-,-\rangle$ is called an inner product on $V$.
\end{definition}

\begin{remark}
Dot product is the standard inner product on $\mathbb{C}^n.$
\end{remark}

\begin{definition}
    Let $(V,\langle-,-\rangle)$ be a complex inner product space.
    \begin{enumerate}
        \item The norm of $x\in V$ is defined as $\norm{x}=\sqrt{\langle x,x\rangle}.$
        \item $x,\,y\in V$ are orthogonal if $\langle x, y\rangle = 0.$
    \end{enumerate}
\end{definition}


\begin{theorem}
    Let $(V,\langle-,-\rangle)$ be a complex inner product space, for $u,v,w\in V$ and $\alpha,\beta\in\mathbb{C}$.
    \begin{enumerate}
        \item Almost linear in 2nd comp: $\langle u,\alpha v+\beta w\rangle=\overline{\alpha}\langle u,v\rangle+\overline{\beta}\langle u,w\rangle$.
        \item $||\alpha u||=|\alpha|||u||$.
        \item $\langle u,v\rangle=0$ for any $v\in V \Leftrightarrow u\equiv 0$.
        \item $\langle u,w\rangle=\langle v,w\rangle$ for any $w\in V \Leftrightarrow u=v$.
    \end{enumerate}
\end{theorem}

\begin{pfbox}
    Left as an exercise.
\end{pfbox}

\begin{theorem}
    Let $(V,\langle-,-\rangle)$ be a complex inner product space, for $u,v\in V$.
    \begin{itemize}
        \item Cauchy-Schwarz: $|\langle u,v\rangle| \le ||u||\cdot||v||$.
        \item Triangle: $||u+v||\le||u||+||v||$.
    \end{itemize}
\end{theorem}

\begin{pfbox}
    For Cauchy-Schwarz, if $v=0,$ then it is trivial. Assume $v\neq 0,$ inspired by the projection vector in $\mathbb{R}^2,$ consider $z=u-\frac{\langle u, v\rangle}{\norm{v}^2}v.$
    First show $z,\,v$ are orthogonal, and by non-negativity, $\langle z,z\rangle\ge 0.$ After some expansion, we have
    $$\norm{u}^2-\frac{\langle u, v\rangle}{\norm{v}^2}\langle v, u\rangle \ge 0,$$
    rearrage and take the square root, then we are done.\\ 

    For triangle inequality,
    \begin{align*}
        ||u+v||^{2} &= \langle u+v,u+v\rangle = \langle u,u\rangle+\langle u,v\rangle+\langle v,u\rangle+\langle v,v\rangle \\
        &= \langle u,u\rangle+2Re(\langle u,v\rangle)+\langle v,v\rangle \\
        &\le ||u||^{2}+2|\langle u,v\rangle|+||v||^{2} \\
        &\le ||u||^{2}+2||u||||v||+||v||^{2} = (||u||+||v||)^{2}.
    \end{align*}
    Hence $||u+v||\le||u||+||v||$.
\end{pfbox}


\begin{theorem}
    Let $(V,\langle-,-\rangle)$ be a complex inner product space. For any $u,v\in V$, we have:
    \begin{itemize}
        \item \textbf{Parallelogram Identity} $\|u+v\|^2 + \|u-v\|^2 = 2(\|u\|^2 + \|v\|^2)$.
        \item \textbf{Generalized Pythagorean Theorem} If $u_1, u_2, \dots, u_n$ are mutually orthogonal (i.e., $\langle u_i, u_j\rangle = 0$ for all $i \neq j$), then 
        $$ \left\| \sum_{i=1}^n u_i \right\|^2 = \sum_{i=1}^n \|u_i\|^2. $$
        \item \textbf{Polarization Identity} 
        $$ \langle u,v\rangle = \frac{1}{4} \left( \|u+v\|^2 - \|u-v\|^2 + i\|u+iv\|^2 - i\|u-iv\|^2 \right). $$
    \end{itemize}
\end{theorem}

\begin{pfbox}
    Left as an exercise.
\end{pfbox}

\begin{theorem}[Jordan-von Neumann]
    Let $V$ be a vector space equipped with a norm $\|\cdot\|$. 
    This norm arises from an inner product (i.e., there exists an inner product on $V$ such that $\|v\| = \sqrt{\langle v,v\rangle}$ for all $v\in V$) if and only if the norm satisfies the parallelogram identity:
    $$ \|u+v\|^2 + \|u-v\|^2 = 2(\|u\|^2 + \|v\|^2) \quad \forall u,v\in V. $$
\end{theorem}

\begin{pfbox}
    Beyond the scope of this note.
\end{pfbox}

\begin{definition}[Orthogonal and Orthonormal Sets]
    Let $(V,\langle-,-\rangle)$ be an inner product space. A subset $S = \{\mathbf{v}_1, \mathbf{v}_2, \dots, \mathbf{v}_k\}$ of $V$ is called an \textbf{orthogonal set} if every pair of distinct vectors is orthogonal, i.e.,
    $$ \langle \mathbf{v}_i, \mathbf{v}_j \rangle = 0 \quad \text{for all } i \neq j. $$
    Furthermore, $S$ is called an \textbf{orthonormal set} if it is an orthogonal set and every vector in $S$ is a unit vector. Compactly, we can write the condition for an orthonormal set using the Kronecker delta $\delta_{ij}$:
    $$ \langle \mathbf{v}_i, \mathbf{v}_j \rangle = \delta_{ij} = \begin{cases} 1, & \text{if } i = j \\ 0, & \text{if } i \neq j \end{cases} $$
\end{definition}

\begin{theorem}
    Let $S = \{\mathbf{v}_1, \mathbf{v}_2, \dots, \mathbf{v}_k\}$ be an orthogonal set of non-zero vectors in an inner product space $V$. Then $S$ is linearly independent.
\end{theorem}

\begin{theorem}
Let $(V,\langle-,-\rangle)$ be a finite-dimensional complex inner product space.
If $S=\{u_1,u_2,\cdots,u_k\}$ is an orthogonal basis of $V$, then for any $v\in V$,
$$v=\sum_{i=1}^{k}\alpha_i u_i, \quad \text{ where } \alpha_i=\frac{\langle v, u_i\rangle}{\norm{u_i}^2}.$$
In particular, if S is an ONB, then $\alpha_i=\langle v, u_i\rangle.$
\end{theorem}

\begin{theorem}
    Let $(V,\langle-,-\rangle)$ be a finite-dimensional complex inner product space. Let $\beta = \{\mathbf{u}_1, \mathbf{u}_2, \dots, \mathbf{u}_n\}$ be an orthonormal basis (ONB) for $V$. For any $\mathbf{v}, \mathbf{w} \in V$, let $[\mathbf{v}]_\beta$ and $[\mathbf{w}]_\beta$ be their coordinate vectors relative to $\beta$. Then the inner product of $\mathbf{v}$ and $\mathbf{w}$ is equal to the standard dot product of their coordinate vectors:
    $$ \langle \mathbf{v}, \mathbf{w} \rangle_V = \langle [\mathbf{v}]_\beta, [\mathbf{w}]_\beta \rangle_{\mathrm{dot}} $$
    \label{coordinate dot product}
\end{theorem}

\begin{theorem}[Gram-Schmidt Process]
    Let $(V,\langle-,-\rangle)$ be an inner product space. Given a linearly independent set $S=\{v_1,v_2,\cdots,v_l\}$ in $V$, if
    \begin{enumerate}
        \item $w_1=v_1$,
        \item $w_k=v_k-\sum_{j=1}^{k-1}\frac{\langle v_k, w_j\rangle}{\norm{w_j}^2}w_j\,$ for $\,2\le k\le l.$
    \end{enumerate} 
    Then:
    \begin{enumerate}
        \item $S'=\{w_1,w_2,\cdots,w_k\}$ is an orthogonal set of non-zero vectors.
        \item $\mathrm{Span}(S)=\mathrm{Span}(S').$
    \end{enumerate}
\end{theorem}

\begin{pfbox}
    We shall prove by induction on $l=|S|.$
    For $l=1,$ there is nothing to prove. Suppose the claim is true for $l=k-1,$ then for $l=k,$
    we need to check:
    \begin{enumerate}
        \item $w_k\neq 0.$
        \item $\langle w_k, w_i\rangle=0$ for $1\le i\le l$.
        \item Take 
        $$w_k\in\mathrm{Span}(\{w_1,w_2,\cdots,w_{k-1},v_{k}\})=\mathrm{Span}(\{v_1,v_2,\cdots,v_{k-1},v_{k}\}),$$ 
        so $\mathrm{Span}(S')\subseteq\mathrm{Span}(S).$
        Note that $\mathrm{Span}(S)$ is given to be linearly independent and $\mathrm{Span}(S')$ is linearly independent by orthogonality of non-zero vectors, so both sets have the same dimension. Hence $\mathrm{Span}(S')=\mathrm{Span}(S).$
    \end{enumerate}
\end{pfbox}

\begin{definition}[Orthogonal Projection]
    Let $(V,\langle-,-\rangle)$ be a finite-dimensional inner product space and $E$ be a subspace of $V$.
    Let $S=\{u_1,u_2,\cdots,u_k\}$ be an ONB of $E$. The projection operator onto $E$ is defined as:
    \[P_E(v)=\sum_{j=1}^{k}\langle v, u_j\rangle u_j.\] 
\end{definition}

\begin{remark}\
    \begin{enumerate}
        \item The projection operator $P_E:V\to E\subseteq V$ is a linear opeator.
        \item $P_E$ is an idempotent operator, i.e. $P^2_E v = P_E v\quad\forall v\in V$.
    \end{enumerate}
\end{remark}

\begin{theorem}
    Let $(V,\langle-,-\rangle)$ be a finite-dimensional inner product space and $P_E:E\to V$ be the projection onto a subspace $E$.
    Then:
    \begin{enumerate}
        \item $P_E(v)\in E$.
        \item $v-P_E(v)$ is orthogonal to every vector in $E$.
    \end{enumerate}
\end{theorem}

\begin{pfbox}
    For 1, it is clear by definition.\\
    For 2, it suffices to check $\langle v-P_E(v), u_j\rangle=0\quad \forall j=1,2,\cdots,k.$
\end{pfbox}


\begin{theorem}[Minimal Distance Property]
  Let $(V,\langle-,-\rangle)$ be a finite-dimensional inner product space and $E$ be a subspace of $V$.
  Fix any $v\in V$, then:
  \begin{enumerate}
    \item Among all vectors in $E$,
    $$\norm{P_E(v)-v}\le\norm{z-v}\quad \forall z\in E.$$
    \item If the equality holds, then $z=P_E(v).$
  \end{enumerate}
\end{theorem}

\begin{pfbox}
    For 1, observe that $\langle z-P_E(v), v-P_E(v)\rangle=0,$ then we can use Pythoagorean theorem.\\
    For 2, the equality holds means $\norm{z-P_E(v)}=0.$ By non-degeneracy, $z=P_E(v).$
\end{pfbox}


\begin{definition}[Orthogonal Complement]
    Let $(V,\langle-,-\rangle)$ be an inner product space and $E$ be a subset of $V$.
    The orthogonal complement of $E$ is defined as:
    \[E^\perp=\{v\in V: \langle v,u\rangle\}=0\quad\forall u\in E.\]
\end{definition}

\begin{remark}\
    \begin{enumerate}
        \item $E$ is not neccessarily a subspace.
        \item $E^\perp$ is a subspace of $V$.
    \end{enumerate}
\end{remark}

\begin{theorem}
 Let $(V,\langle-,-\rangle)$ be a finite-dimensional inner product space and $E$ be a subspace of $V$.
 \begin{enumerate}
    \item $v-P_E(v)\in E^\perp.$
    \item $\Ker(P_E)=E^\perp.$
\end{enumerate}
\end{theorem}

\begin{pfbox}
    For 1, it is followed by $v-P_E(v)$ is orthogonal to any vector in $E$.\\
    For 2, take $v\in \Ker(P_E),$ then $v-P_E(v)=v\in E^\perp.$
    Conversely, take $v\in E^\perp,$ then $P_E(v)=\underbrace{v-P_E(v)}_{\in E^\perp}+v=P_E(v)\in E^\perp.$\\
    By non-degeneracy, $P_E(v)=0.$
\end{pfbox}


\begin{theorem}[Orthogonal Decomposition]
 Let $(V,\langle-,-\rangle)$ be a finite-dimensional inner product space and $E$ be a subspace of $V$.
 \begin{enumerate}
    \item $V=E\oplus E^\perp$
    \item $\dim(V)=\dim(E)+\dim(E^\perp)..$
\end{enumerate}
\end{theorem}

\begin{pfbox}
    For 1, observe that $v=P_E(v)+(v-P_E(v))$ for any $v\in V.$ To show $E$ and $e=E^\perp$ has trivial intersection, we take $v\in E\cap E^\perp,$ then $\langle v,v\rangle=0.$ By non-degeneracy, $v=0.$\\
    For 2, it follows from 1 directly.
\end{pfbox}



\begin{theorem} 
    Let $A\in M_{m\times n}(\mathbb{R}),\,x\in\mathbb{R}^n\,y\in\mathbb{R}^m.$
    \begin{enumerate}
        \item $\langle Ax, y\rangle_{\mathrm{dot}} = \langle x, A^\top y\rangle_{\mathrm{dot}}.$
        \item $\rank(A^\top A)=\rank(A).$
        \item If $rank(A)=n$, then $A^\top A$ is inveritble.
    \end{enumerate}
\end{theorem}

\begin{pfbox}
    For 1 and 3, it is an easy exercise.\\
    For 2, we first show $\ker(A^\top A)=\ker(A).$\\
    For any $v\in\ker(A^\top A),\,\langle A^\top Av, v\rangle=0.$ By 1, $\langle Av, Av\rangle=0$, so $Av=0$ by degeneracy. The converse is trivial.
    In addition, both $A$ and $A^\top A$ shares the same dimension of domain (The matrics have the same number of columns). By rank-nullity, $\rank(A^\top A)=\rank(A).$
\end{pfbox}


\begin{question}[Best Solution Problem]
    Given a system of equations $X\beta=y$ to be inconsistent, where $X\in M_{m\times n}(\mathbb{R}),\, y\in \mathbb{R}^m$. Find $\beta_0\in\mathbb{R}^n$ that minimizes $\norm{X\beta-y},$ i.e.
    $$\norm{X\beta_0-y}\le\norm{X\beta-y}$$ 
    for any $\beta\in\mathbb{R}^n.$
\end{question}

\begin{remark}
    The norm is induced by the usual dot product on real vector space.
\end{remark}

To solve the best solution problem, we employ the following method:
\begin{theorem}[Least Square Methods]
    Let $X\in M_{m\times n}(\mathbb{R}),\, y\in\mathbb{R}^m.$ Suppose $\beta=\beta_0\in\mathbb{R}^n$ is a minimizer of $\norm{y-X\beta},$ then:
    \begin{enumerate}
        \item $\beta_0$ satisfies $X^\top X\beta=X^\top y;$
        \item If $rank(X)=n$, then $\beta_0=(X^\top X)^{-1}X^\top y.$
    \end{enumerate}
\end{theorem}

\begin{pfbox}
    For 1, let $W=\Ima(X)=\{X\beta:\beta\in\mathbb{R}^n\}.$
    By the minimal distance property, 
    \begin{align*}
                &\norm{P_W(y)-y}\le\norm{w-y}\,\forall w\in W\\
     \Rightarrow &\norm{P_W(y)-y}\le\norm{X\beta-y}\,\forall \beta\in\mathbb{R}^n.
    \end{align*}
    So it suffices to find $\beta_0$ such that $P_W(y)=X\beta_0.$
    This is possible since $P_W(y)\in W.$
    Let $P_W(y)=X\beta_0$ for some $\beta_0\in\mathbb{R}^n.$
    Note that $$\langle P_W(y)-y, w\rangle = \langle X\beta_0-y, X\beta\rangle = \langle X^\top (X\beta_0-y), \beta\rangle=0\quad\forall \beta\in\mathbb{R}^n,$$
    then we have $X^\top (X\beta_0-y)=0.$\\

    For 2, it is an easy exercise.
\end{pfbox}



\begin{question}[Minimal Solution Problem]
        Given a system of equations $Ax=b$ to be infinitely many solutions, where $A\in M_{m\times n}(\mathbb{R}),\, b\in \mathbb{R}^m$. Find $u\in\mathbb{R}^n$ that s.t. $\norm{u}\le\norm{x}$ for all solutions $x$ to the system.
\end{question}

\begin{theorem}
    Let $A\in M_{m\times n}(\mathbb{R}),$ then $\Ker(A)=(\Ima(A^\top))^\perp.$ 
\end{theorem}

\begin{pfbox}
    Left as an exercise.
\end{pfbox}


\begin{theorem}
  Given a system of equations $Ax=b$ to be infinitely many solutions, where $A\in M_{m\times n}(\mathbb{R}),\, b\in \mathbb{R}^m$. 
  \begin{enumerate}
    \item There exists a unique minimal solution $u$ to the system and $u\in\Ima(A^\top).$
    \item Let $v$ be a solution to $AA^\top v=b,$ then $A^\top v=u.$ 
  \end{enumerate}
\end{theorem}

\begin{pfbox}
    For 1, for any solution $x_0$ to the system, we can write a general solution $u=x_0-v$ for some $v\in\Ker(A).$
    It suffices to find $v$ such that $\norm{x_0-v}\le\norm{x_0-w}$ for all $w\in\Ker(A).$ Consider the projection $v=P_{\Ker(A)}(x_0).$
    By the minimal distance property, $v$ is unique, so $u=x_0-v\in(\Ker(A))^\perp=\Ima(A^\top)$ is also unique.\\

    For 2, suppose $AA^\top v=b$ and $Au=b$ by construction, then $A^\top v-u\in\Ker(A)\cap\Ima(A^\top)=\{0\}.$
    Hence $A^\top v=u.$ 
\end{pfbox}


\begin{remark}\
Why $AA^\top v=b$ is consistent?
Note that $\Ker(A^\top)=\Ker(AA^\top).$ 
Use $\Ker(A)=\Ima(A^\top)^\perp$ and take orthogonal complement on both sides, one has:
$\Ima(AA^\top)=\Ima(A),$ then we are done.
\end{remark}

\begin{theorem}
    Let $A\in M_n(\mathbb{C})$ and $v,w\in\mathbb{C}^n,$ then
    $$\langle Av,w\rangle = \langle v, \overline{A^\top}w\rangle.$$
\end{theorem}

\begin{definition}[Adjoint Operator]
    Let $A\in M_n(\mathbb{C}).$ The adjoint of $A$ is defined by its conjugate transpose,
    denoted by $A^*=\overline{A^\top}.$

    In general, the adjoint of a linear operator $T$ is an operator $T^*:V\to V$ if
    \[\langle Tv,w\rangle = \langle v, T^*w\rangle\,\,\forall\,\, v,w\in V\]
\end{definition}

\begin{theorem}
    For each linear operator, its adjoint operator exists and is unique.
\end{theorem}

\begin{pfbox}
    For existence, we first fix an ONB $\beta,$ and define $[T^*]^\beta_{\beta} = \overline{([T]^\beta_{\beta})^\top}$.
    Use theorem~\eqref{coordinate dot product}, we check that 
    $$\langle Tv,w\rangle = \langle [Tv]_\beta,[w]_\beta\rangle = \langle [v]_\beta, [T^*]^\beta_{\beta}[w]_\beta\rangle = \langle v, T^*w\rangle.$$

    For uniqueness, we left as an exercise.
\end{pfbox}


\begin{theorem}
    Let $T,\,U:V\to V$ be two linear operators on an inner product space $V$. We have:
    \begin{enumerate}
        \item $(T+U)^* = T^* + U^*.$
        \item $(cT)^*=\bar{c}T^*$ for any $c\in\mathbb{C}.$
        \item $(TU)^* = U^*T^*.$
        \item $(T^*)^* = T.$
    \end{enumerate}
\end{theorem}

\begin{theorem}
   Let $S$ be an ONB of $V$, and $S'$ be another basis of $V$. TFAE:
   \begin{enumerate}
    \item $S'$ is also an ONB of $V$.
    \item The change of coordinate matrix $U$ from $S'$ to $S$ satisfies $U^*U=U^*U=I_n$ (i.e. $U$ is unitary).
   \end{enumerate}
\end{theorem}

\begin{definition}
    Let $A\in M_n(\mathbb{C}),$
    \begin{enumerate}
        \item $A$ is normal if $AA^*=A^*A$.
        \item $A$ is unitary if $AA^*=A^*A = I_n$ (i.e. $A^*=A^{-1}$).
        \item $A$ is Hermitian/self-adjoint if $A=A^*$. 
    \end{enumerate}
\end{definition}

\begin{theorem}
    For a finite-dimensional inner product space $V$, the set of unitary operators forms a group.
    That is,
    $$U(V)=\{T:V\to V: T\text{ is unitary}\} \text{ is}$$
    \begin{enumerate}
        \item closed under composition,
        \item closed under inversion.
    \end{enumerate}
\end{theorem}

\begin{theorem}
    Let $T\in U(V)$ and $\lambda\in\mathrm{Spec}(T).$ Then $|\lambda|=1,$ i.e. 
    $$\lambda=e^{i\theta}=\cos\theta + i\sin\theta\,\, \forall\theta\in\mathbb{R}.$$
\end{theorem}

\begin{pfbox}
Compute $\langle Tv, Tv\rangle = \langle v, T^*Tv\rangle =\langle v,v\rangle$ by unitarity of $T$.
On the other hand, $\langle Tv, Tv\rangle = \langle \lambda v, \lambda v\rangle = \lambda\overline{\lambda} \norm{v}.$
Compare them and $v\neq 0,$ we obtain $|\lambda| = 1.$
\end{pfbox}

\begin{lemma}
Let $T$ be a linear operator on a finite-dimensional complex inner product space
$(V,\langle -, -\rangle)$. If $\langle Tv,Tw\rangle = \langle v,w\rangle
\,\,\forall v,w\in V$, then $T$ is unitary.
\end{lemma}

\begin{pfbox}
Suppose $\langle Tv,Tw\rangle = \langle v,w\rangle
\quad \forall v,w\in V$.
Then
\begin{align*}
& \langle v,T^*Tw\rangle = \langle v,w\rangle\\
& \langle v,(T^*T-I)w\rangle=0 \quad \forall v,w\in V.
\end{align*}

Since this holds for every $v\in V$,
\begin{align*}
    &(T^*T-I)w=0 \quad \forall w\in V \\ 
    & T^*T=I.
\end{align*}
\end{pfbox}




\begin{theorem}
    Let $(V,\langle-,-\rangle)$ be a finite-dimensional inner product space and $T: V \to V$ be a linear operator. 
    Then $T$ is a unitary operator if and only if $T$ is an isometry. That is,
    $$ T \in U(V) \iff \norm{Tv} = \norm{v} \quad \forall v \in V. $$
\end{theorem}

\begin{pfbox}
($\Rightarrow$) Suppose $T$ is unitary, $T^*T=I$. Then for any $v\in V$,
\begin{align*}
\|Tv\|^2 =\langle Tv,Tv\rangle =\langle v,T^*Tv\rangle =\langle v,v\rangle =\|v\|^2.
\end{align*}\\
($\Leftarrow$)
Assume
\[
\|Tv\|=\|v\|
\qquad \forall v\in V.
\]

For any $x,y\in V$, by the polarization identity,
\begin{align*}
\langle x,y\rangle
&=
\frac14\Big(
\|x+y\|^2-\|x-y\|^2
\Big)  \\
&\quad
+\frac{i}{4}
\Big(
\|x+iy\|^2-\|x-iy\|^2
\Big).
\end{align*}

Using the hypothesis,
\begin{align*}
\langle x,y\rangle
&=
\frac14
\Big(
\|T(x+y)\|^2
-
\|T(x-y)\|^2
\Big) \\
&\quad
+\frac{i}{4}
\Big(
\|T(x+iy)\|^2
-
\|T(x-iy)\|^2
\Big).
\end{align*}

Since $T$ is linear,
\begin{align*}
\langle x,y\rangle
&=
\frac14
\Big(
\|Tx+Ty\|^2
-
\|Tx-Ty\|^2
\Big) \\
&\quad
+\frac{i}{4}
\Big(
\|Tx+iTy\|^2
-
\|Tx-iTy\|^2
\Big) \\
&=
\langle Tx,Ty\rangle.
\end{align*}

Therefore
\[
\langle Tx,Ty\rangle
=
\langle x,y\rangle
\qquad \forall x,y\in V.
\]

By the previous lemma, $T$ is unitary.
\end{pfbox}

\begin{theorem}
Let $A\in M_{n}(\mathbb{C})$. The following are equivalent.
\begin{enumerate}
    \item $AA^{*}=A^{*}A=I_{n}.$
    \item The rows of A form an orthonormal basis for $\mathbb{C}^{n}$.
    \item The columns of A form an orthonormal basis for $\mathbb{C}^{n}$.
\end{enumerate}
\end{theorem}

\begin{pfbox}
Let $A = \begin{pmatrix} \text{---} & r_1 & \text{---} \\ \text{---} & r_2 & \text{---} \\ & \vdots & \\ \text{---} & r_n & \text{---} \end{pmatrix}$.
Observe that:
\[
AA^{*} = \begin{pmatrix} \langle r_1, r_1 \rangle_{\text{dot}} & \langle r_1, r_2 \rangle_{\text{dot}} & \dots \\ \vdots & \ddots & \vdots \\ \dots & \dots & \langle r_n, r_n \rangle_{\text{dot}} \end{pmatrix}.
\]
(1) $\Leftrightarrow$ (2): $AA^{*}=I_{n} \Leftrightarrow \langle r_i, r_j \rangle_{\text{dot}} = \delta_{ij} \Leftrightarrow \{r_1, r_2, \dots, r_n\}$ forms an ONB.
Similarly, replace rows by columns and $AA^*$ by $A^*A$ in the above argument to show (1) $\Leftrightarrow$ (3).
\end{pfbox}

\begin{theorem}
If $T$ is a Hermitian operator on $V$, then all eigenvalues of $T$ are real.
\end{theorem}

\begin{pfbox}
Let $T$ be a Hermitian operator. Take any eigenvalue $\lambda \in \Spec(T)$, so $Tv = \lambda v$ for some $v \neq 0$.
Consider:

\[
\langle Tv, v \rangle = \langle \lambda v, v \rangle = \lambda \langle v, v \rangle
\]
On the other hand, since T is Hermitian ($T = T^*$):

\[
\langle Tv, v \rangle = \langle v, T^*v \rangle = \langle v, Tv \rangle = \langle v, \lambda v \rangle = \bar{\lambda} \langle v, v \rangle
\]
Hence, $\lambda \langle v, v \rangle = \bar{\lambda} \langle v, v \rangle$. Since $v \neq 0$, we have $\langle v, v \rangle \neq 0$, which implies:

\[
\lambda = \bar{\lambda} \Leftrightarrow \lambda \in \mathbb{R}.
\]
\end{pfbox}

\begin{theorem}
Let T be a Hermitian operator on V. Then any eigenvectors that correspond to different eigenvalues of T are orthogonal.
\label{orthogonal eigenvectors Hermitian case}
\end{theorem}

\begin{pfbox}
\textbf{WTS:} If $Tu = \lambda u$ and $Tw = \mu w$, where $\lambda \neq \mu$, then $\langle u, w \rangle = 0$.
Consider:
\[
\langle Tu, w \rangle = \langle \lambda u, w \rangle = \lambda \langle u, w \rangle
\]
Also:
\[
\langle u, T^*w \rangle = \langle u, Tw \rangle = \langle u, \mu w \rangle = \bar{\mu} \langle u, w \rangle = \mu \langle u, w \rangle
\]
(Since T is Hermitian, its eigenvalues are real, so $\bar{\mu} = \mu \in \mathbb{R}$).

Thus, $\lambda \langle u, w \rangle = \mu \langle u, w \rangle$, which implies:

\[
(\lambda - \mu) \langle u, w \rangle = 0
\]
Since $\lambda \neq \mu$, we must have $\langle u, w \rangle = 0$.
\end{pfbox}


\begin{proposition}
Let $(V,\langle-,-\rangle)$ be a complex finite-dimensional inner product space and $T:V\to V$ be a normal operator. 
Then:
\begin{enumerate}
    \item $\|T(v)\|=\|T^*(v)\| \,\, \forall v\in V.$
    \item If $\lambda$ is an eigenvalue of $T$, then $\overline{\lambda}$ is an eigenvalue of $T^*.$
\end{enumerate}
\label{conjugate transpose of normal matrix}
\end{proposition}


\begin{pfbox}
Fix any $v\in V$. Since $T$ is normal, $TT^*=T^*T$,
\begin{align*}
\|Tv\|^2 &=\langle Tv,Tv\rangle =\langle v,T^*Tv\rangle=\langle v,TT^*v\rangle \\
&=\langle T^*v,T^*v\rangle =\|T^*v\|^2.
\end{align*}

Next, let $\lambda\in \Spec(T)$, then there exists $v\neq 0$ such that $Tv=\lambda v \Rightarrow (T-\lambda I)v=0$.
One can verify that $T-\lambda I$ is normal.\\
By the first part,
$ \|(T-\lambda I)v\| = \|(T-\lambda I)^*v\|$.
Since
$(T-\lambda I)v=0$, we obtain $\|(T-\lambda I)^*v\|=0$.
By non-degeneracy, $(T^*-\overline{\lambda}I)v=0$.
Since $v\neq0$, $\overline{\lambda}\in \Spec(T^*)$.
\end{pfbox}

\begin{theorem}
Let $T$ be a normal operator on V. Then any eigenvectors that correspond to different eigenvalues of T are orthogonal.
\end{theorem}

\begin{pfbox}
We can apply the same argument as in the proof of theorem~\eqref{orthogonal eigenvectors Hermitian case},
combining with the previous proposition.
\end{pfbox}

\begin{remark}
    This is the generalization of theorem~\eqref{orthogonal eigenvectors Hermitian case}.
\end{remark}


\begin{definition}
Let $A\in M_{n}(\mathbb{R})$. Then
\begin{enumerate}
    \item A is called \textbf{orthogonal} if $AA^{t}=A^{t}A=I_{n}$ (orthogonal = real unitary).
    \item A is called \textbf{symmetric} if $A^{t}=A$ (symmetric = real Hermitian).
\end{enumerate}
\end{definition}

\begin{remark}
For convenience, we often write
\[
O(n) \overset{\text{def}}{=} \{A \in M_{n}(\mathbb{R}) : A^{t}A = AA^{t} = I_{n}\}
\]
to be the set of n by n orthogonal matrices. This is also a 'group' (closed under multiplication, inversion).
\end{remark}

\begin{theorem}
If $A\in M_{n}(\mathbb{R})$ is an orthogonal matrix, then
\begin{enumerate}
    \item all the eigenvalues $\lambda$ of A satisfy $|\lambda|=1$;
    \item $\det(A) = \pm 1$;
    \item the columns/rows of A form an orthonormal basis of $\mathbb{R}^{n}$.
\end{enumerate}
\end{theorem}

\begin{remark}
Although A is real, it does not imply its eigenvalues are real (e.g., rotation matrix).
\end{remark}

\begin{pfbox}
\textbf{Proof of 1 \& 3:} They follow directly from properties of unitary matrices. \\
\textbf{Proof of 2:} By def, if A is orthogonal, $AA^{t}=I_{n}$.
\begin{align*}
&\det(AA^{t}) = \det(I_{n})\\ \Rightarrow 
&\det(A)\det(A^{t}) = 1\\ \Rightarrow 
&\det(A) = \pm 1.
\end{align*}
\end{pfbox}

\begin{theorem}
If $A\in M_{n}(\mathbb{R})$ is a symmetric matrix, then
\begin{enumerate}
    \item all the eigenvalues of A are real;
    \item any eigenvectors that correspond to different eigenvalues of A are orthogonal (with respect to the real dot product).
\end{enumerate}

\end{theorem}

\begin{pfbox}
Symmetric = real Hermitian. (As A is real, $A^* = \overline{A^t} = A^t = A$).
These matrices enjoy all properties of Hermitian matrices.
\end{pfbox}



\subsection{Spectral Theorem}

\begin{definition}\ \vspace{-2em}
\begin{enumerate}
    \item $A\in M_{n}(\mathbb{C})$ is called \textbf{unitarily diagonalizable} if there exists an unitary matrix U ($U^* = U^{-1}$) such that $U^{*}AU$ is a diagonal matrix.
    \item $A\in M_{n}(\mathbb{R})$ is called \textbf{orthogonally diagonalizable} if there exists an orthogonal matrix P ($P^t = P^{-1}$) such that $P^{t}AP$ is a diagonal matrix.
\end{enumerate}
\end{definition}



\begin{theorem}[Refined Schur Triangulation]
Let $T: V \rightarrow V$ be a linear operator. Then there exists an orthonormal basis $\beta$ of V such that the matrix $[T]_{\beta}^{\beta}$ is an upper triangular matrix.
\label{refined schur}
\end{theorem}

\begin{pfbox}
By Schur Triangulation,
$\exists$ a basis $\beta' = \{v_1, v_2, \dots, v_n\}$ of V s.t.\\
$[T]_{\beta'}^{\beta'}$ is upper triangular.
Apply Gram-Schmidt and normalize,\\
we obtain an ONB $\beta=\{u_1, u_2, \dots, u_n\}$.\\
We want to show $T(u_k) \in \spn\{u_1, u_2, \dots, u_k\}$ for each $k=1,2,\dots,n$.
We proceed by induction on $k$. $k=1$ is trivial.
Now assume the claim is true up to $k-1$,
by Gram-Schmidt, $u_k-v_k \in \spn\{u_1, u_2, \dots, u_{k-1}\}$.
So $T(u_k)-T(v_k) \in \spn\{T(u_1), T(u_2), \dots, T(u_{k-1})\} \subseteq \spn\{u_1, u_2, \dots, u_{k-1}\}$ by induction hypothesis.
Note that $T(v_k) \in \spn\{v_1, v_2, \dots, v_k\} = \spn\{u_1, u_2, \dots, u_k\}$.
Hence, $T(u_k) \in \spn\{u_1, u_2, \dots, u_k\}$ for each $k=1,2,\dots,n$.
\end{pfbox}


\begin{theorem}[Hermitian Spectral Theorem]
If $A\in M_{n}(\mathbb{C})$ is Hermitian, then $A$ is unitarily diagonalizable.
\end{theorem}

\begin{pfbox}
By refined Schur's Triangulation, $\exists U \in U(n)$ s.t. $U^{*}AU = D$ is upper triangular.
We should observe that $D$ is also Hermitian, and hence diagonal.
\end{pfbox}

\begin{theorem}
Let $A\in M_{n}(\mathbb{C})$ be an upper triangular matrix. Then A is normal if and only if A is diagonal.
\label{A is upper trig. then A is normal iff diagonal}
\end{theorem}

\begin{pfbox}
Let A be upper triangular.\\
$(\Leftarrow)$ Trivial.\\
$(\Rightarrow)$ Suppose A is normal as well, $A^{*}A = AA^{*}$. Let $A=(a_{ij})$.
Compare the $(1,1)$ entries of $AA^{*}$ and $A^{*}A$:
\[
(AA^{*})_{11} = |a_{11}|^2 + |a_{12}|^2 + \dots + |a_{1n}|^2
\]

\[
(A^{*}A)_{11} = |a_{11}|^2
\]
As A is normal, $AA^{*} = A^{*}A$, which forces:
\[
|a_{12}|^2 + |a_{13}|^2 + \dots + |a_{1n}|^2 = 0 \Rightarrow a_{1j} = 0 \text{ for } j \geq 2.
\]
Repeat this argument to other diagonal entries of $AA^{*}$ and $A^{*}A$. We can deduce $a_{ij} = 0$ whenever $i \neq j$.
So $A$ is diagonal.
\end{pfbox}


\begin{theorem}[$\mathbb{C}$omplex Spectral Theorem]
$A\in M_{n}(\mathbb{C})$ is unitarily diagonalizable if and only if A is normal.
\end{theorem}

\begin{remark}
Equivalently, in terms of 'operator language', the Complex Spectral Theorem states that $T:V\rightarrow V$ is a normal operator if and only if there exists an orthonormal basis $\beta$ such that $[T]_{\beta}^{\beta}$ is diagonal.
\end{remark}


\begin{pfbox}
($\Rightarrow$): Let A be normal. By refined Schur triang., $\exists U \in U(n)$ s.t. $U^{*}AU = D$ is upper triangular.
Observe that $D$ is normal, by theorem~\eqref{A is upper trig. then A is normal iff diagonal} $D$ must be diagonal.

($\Leftarrow$): Exercise by direct compute.
\end{pfbox}


\begin{theorem}[$\mathbb{R}$eal Spectral Theorem]
Let $A\in M_{n}(\mathbb{R})$. The following are equivalent.
\begin{enumerate}
    \item $A\in M_{n}(\mathbb{R})$ is symmetric. (i.e. $A^{t}=A$)
    \item There exists $P\in O(n)$ such that $P^{t}AP$ is diagonal (or so-called 'orthogonally diagonalizable').
\end{enumerate}

\end{theorem}

\begin{pfbox}
(1) $\Rightarrow$ (2): Let A be symmetric.
$A^{*} = \overline{A^{t}} = A^{t} = A$. So A is Hermitian (hence normal).
By Complex Spectral Theorem, $\exists U \in U(n)$ such that $U^{*}AU = D$ is diagonal.
Recall the columns of U consist of eigenvectors of A.
As A is symmetric, all its eigenvalues are real. As $A \in M_{n}(\mathbb{R})$, the eigenvectors (solutions to $Ax = \lambda x$) can be chosen to be real.
Hence, we can choose U to be a "real" matrix.
Thus, U is orthogonal (real unitary), $U \in O(n)$.\\

(2) $\Rightarrow$ (1): Suppose A is orthogonally diagonalizable. 
$\exists P \in O(n)$ such that $P^{t}AP = D$ is diagonal, where $P^{t} = P^{-1}$.
Note that $A = PDP^{t},\, A^{t} = (PDP^{t})^{t} = (P^{t})^{t}D^{t}P^{t} = PDP^{t} = A$ since $D$ is diagonal.
A is symmetric.
\end{pfbox}

\begin{proposition}
If $A\in M_{n}(\mathbb{C})$ is normal, then A and $A^{*}$ share the same eigenvectors.
\end{proposition}

\begin{pfbox}
Given A is normal, by Spectral Theorem, $\exists\, U$ unitary such that $U^{*}AU = D$ is diagonal.
Since the columns of $U$ forms an eigen-basis (and ONB as well) of $A$
by taking the conjugate transpose on both sides, we have $U^{*}A^{*}U = \overline{D}$,
$A^{*}$ can be diagonalized by the same matrix $U$. 
Hence, the columns of $U$ are eigenvectors of $A^{*}$ as well.

Explicitly,  fix every eigenvector $v$ of $A$ with eigenvalues $\lambda$, write
$Av = \lambda v$, since $(A-\lambda I)$ is also normal, by proposition~\eqref{conjugate transpose of normal matrix}, $\norm{(A-\lambda I)v} = \norm{(A-\lambda I)^*v} = 0 \Rightarrow A^*v = \bar{\lambda}v$.
\end{pfbox}



\begin{theorem}[Riesz Representation (Finite-Dimensional Version)]
Let $(V,\langle-,-\rangle)$ be a finite-dimensional complex inner product space. For each linear functional $\varphi\in V^{*}$, there exists a unique $u\in V$ such that
\[
\varphi(v)=\langle v,u\rangle \quad \text{for all } v\in V.
\]
\end{theorem}

\begin{remark}
In other words, there is a natural identification between the dual space $V^{*}$ and V.
\end{remark}

\begin{pfbox}
\textbf{Existence} Fix $\{v_1, v_2, \dots, v_n\}$ to be an ONB of V.
Then for any $v \in V$:
\begin{align*}
&v = \sum_{i=1}^{n}\langle v,v_i\rangle v_i\\\Rightarrow 
&\varphi(v) = \varphi\left(\sum_{i=1}^{n}\langle v,v_i\rangle v_i\right) = \sum_{i=1}^{n}\langle v,v_i\rangle \varphi(v_i)
\end{align*}

Let $u = \sum_{i=1}^{n}\overline{\varphi(v_i)}v_i$.
Then $\langle v, u \rangle = \left\langle v, \sum_{i=1}^{n}\overline{\varphi(v_i)}v_i \right\rangle = \sum_{i=1}^{n} \varphi(v_i)\langle v, v_i \rangle = \varphi(v)$.
We are done by taking u.

\textbf{Uniqueness} Suppose $\exists\, u_1, u_2$ s.t. for all $v \in V$:
\begin{align*}
&\varphi(v) = \langle v, u_1 \rangle = \langle v, u_2 \rangle\\ \Rightarrow
&\langle v, u_1 - u_2 \rangle = 0 \quad \forall v \in V  
\end{align*}

Take $v = u_1 - u_2$, by non-degeneracy, we obtain $u_1 = u_2$.
\end{pfbox}
